Large language model agents are increasingly used for retrieval-augmented question answering, but most Agentic RAG systems are designed around text-centric retrieval: fetch a passage, read it, and continue. Exploratory question answering over tabular data lakes behaves differently. A useful table may require multiple turns of schema inspection, file-handle resolution, SQL or Python execution, parsing repair, and evidence validation before it can contribute to an answer. We argue that this creates a stopping-rule problem: even when strong retrieval tools expose the right sources, agents waste tool calls switching files, rechecking evidence, over-searching, or debugging tabular operations without knowing when a source interaction is complete.

We introduce \textbf{Tabular MOMO}, a runtime-control framework for exploratory QA over data lakes. MOMO separates planning from data interaction through bounded subagent contracts. A planner decomposes the user question and delegates source discovery to search workers and extraction to inspection workers. Search workers are limited to retrieval and lightweight peeking, while inspection workers are restricted to explicit source-family identifiers and exact file references. Each subagent operates under a hard tool budget and returns a compact structured payload containing answer fragments, missing outputs, evidence, retry recommendations, and execution statistics. Raw schemas, SQL errors, row dumps, and parsing loops remain isolated from the planner, forming a \emph{context firewall}.

MOMO is implemented as a modular extension over an existing DataLakeAgent evaluation harness via five generic extension hooks, leaving the baseline agent unmodified. It exposes profile-enriched file peeking, structured tool-use records, source-commitment reflections, and multi-axis ablations over search, planning, computation, and runtime-control settings. We position tabular QA as an execution-discipline problem, not merely a retrieval problem, and present MOMO as an implementation-backed architecture for making exploratory data-lake agents more bounded, auditable, and failure-aware.
